\vspace{-2mm}
\section{Conclusions}\label{sec:conclusion}

% Existing GraphRAG solutions either suffer from limited accuracy and applicability or incur high computation costs, as they adopt a single fixed graph traversal method for all graph questions  or invoke the LLM at every traversal step. To address their limitations, we systematically analyze the graph questions and propose a taxonomy that categorizes the graph questions into four distinct types, based on both query target (i.e., entity or relation) and contents (i.e., specific or abstract). Building on this taxonomy, we design \name{}, a solution that dynamically selects the most suitable graph traversal method for each question type and executes the graph traversal operations with a unified interface and backend. Extensive experiments show that \name{} consistently achieves higher generation quality than SOTA  GraphRAG methods while maintaining low end-to-end latency and token consumption across all four question types.

% \name{} is essentially a query planner for GraphRAG, which is extensively studied in database but less explored for RAG and GraphRAG. We envision that holistic query categorization and planning may be promising directions for processing diverse and complex RAG and GraphRAG queries. 

% Existing GraphRAG solutions are primarily evaluated on current KGQA benchmarks which suffer from limited diversity of graph questions. This hinders a comprehensive evaluation of different GraphRAG methods and leads to biased design of GraphRAG solutions. To this end, we systematically investigate the underlying question patterns from real-world historical questions and propose four basic question patterns that can express or compose any valid graph questions.
In this work, we develop a new benchmark, \benchmarkname{}, for general GraphRAG workloads, covering a wide range of graph question patterns to enable more reliable evaluation of GraphRAG methods. Furthermore, we propose a novel GraphRAG solution, \name{}, which excels at providing high-quality responses to diverse GraphRAG questions. \name{} is an adaptive framework that automatically identifies the question pattern and dynamically prompts the LLM to generate appropriate Cypher queries to retrieve the relevant context from the knowledge graphs. Experiment results show that \name{} achieves high response quality with a low response latency and token usage.
